\section{可完整积分的平面界面损失谱}
\label{chap:feqo-eels}

令真空占据 \(x>0\)，均匀各向同性材料占据 \(x<0\)。电子以速度 \(v\hat{\vb{z}}\) 沿 \(x=b>0,y=0\) 直线运动。材料的相对介电函数为 \(\varepsilon(\omega)\)。先研究主导波数 \(\kappa\gg\omega/c\) 的非迟滞近场：电场可写成 \(-\nabla\varphi\)，但仍保留材料随频率变化的 \(\varepsilon\)。这个近似要求主要贡献来自 \(b^{-1}\) 或 \(\omega/v\) 所选波数远大于 \(\omega/c\)；若两者接近，必须使用完整 Maxwell Green 张量。

把频域电势沿界面 Fourier 展开为 \(\varphi(x,y,z,\omega)=\int(\dd k_y\dd k_z/(2\pi)^2)e^{\ii(k_yy+k_zz)}\widetilde\varphi(x,\vb{k}_\parallel,\omega)\)。\cref{eq:feqo-electron-current} 对应的电荷密度为 \(\rho=q\delta(x-b)\delta(y)e^{\ii\omega z/v}/v\)。故源的 \(z\) 向 Fourier 变换含 \(2\pi\delta(k_z-\omega/v)\)：电子轨迹直接固定 \(k_z=\omega/v\)。令 \(\kappa=\sqrt{k_y^2+k_z^2}\)，真空半空间的 Poisson 方程为
\[
(\partial_x^2-\kappa^2)\widetilde\varphi=-\widetilde\rho/\varepsilon_0.
\]
在 \(x=0\) 处要求电势连续、法向位移 \(\varepsilon\partial_x\varphi\) 连续，解的散射部分为
\begin{equation}
\widetilde\varphi_{\rm ind}(x,k_y,k_z,\omega)
=\frac{q}{2\varepsilon_0v\kappa}
\frac{1-\varepsilon(\omega)}{1+\varepsilon(\omega)}
e^{-\kappa(x+b)}\,2\pi\delta(k_z-\omega/v),\quad x>0.
\label{eq:feqo-planar-image-potential}
\end{equation}
分母 \(2\kappa\) 来自一维方程的点源导数跳跃；分式是电势的静电反射系数，记为 \(r_{\rm es}\)。若 \(\varepsilon=1\)，材料与真空无区别，诱导势为零。

反射系数值得从边界条件亲手复核。在源与界面之间，源本身的 Fourier 势与它激发的反射势分别正比于 \(e^{\kappa x}\) 与 \(r_{\rm es}e^{-\kappa x}\)；材料一侧的衰减解正比于 \(t_{\rm es}e^{\kappa x}\)。在 \(x=0\) 连续电势得到 \(1+r_{\rm es}=t_{\rm es}\)。令界面无自由表面电荷，法向位移连续给 \(1-r_{\rm es}=\varepsilon t_{\rm es}\)。消去 \(t_{\rm es}\) 才得到 \(r_{\rm es}=(1-\varepsilon)/(1+\varepsilon)\)。在第二个条件中把导数的负号漏掉，会把表面模的条件 \(1+\varepsilon=0\) 错写成别的式子。

由 \(E_z=-\partial_z\varphi_{\rm ind}\) 得 \(E_z=-\ii(\omega/v)\varphi_{\rm ind}\)。把它放入\cref{eq:feqo-eels-work}，对电子经过的长度 \(L\) 积分，再除以 \(L\hbar\omega\)，得到单位轨迹长度、单位角频率的单量子损失概率
\begin{equation}
\frac{\dd^2\Gamma_{\rm EELS}}{\dd z\,\dd\omega}
=\frac{q^2}{2\pi\varepsilon_0\hbar v^2}
\int_{-\infty}^{\infty}\frac{\dd k_y}{2\pi}
\frac{e^{-2b\sqrt{k_y^2+(\omega/v)^2}}}{\sqrt{k_y^2+(\omega/v)^2}}
\bigl[-\operatorname{Im}r_{\rm es}(\omega)\bigr].
\label{eq:feqo-planar-eels}
\end{equation}
对被动材料 \(\operatorname{Im}\varepsilon\geq0\)，直接算分式的虚部给 \(-\operatorname{Im}r_{\rm es}=2\operatorname{Im}\varepsilon/|1+\varepsilon|^2\geq0\)，所以本例的概率非负。这里使用的是由两个边界条件求出的具体散射核；它的正性不能从任意散射 Green 张量单独正定来猜。

还可将横向波数积分做到底。令 \(k_y=(\omega/v)\sinh u\)，则 \(\dd k_y/\sqrt{k_y^2+(\omega/v)^2}=\dd u\)。定义修正 Bessel 函数 \(K_0(x)=\int_0^\infty e^{-x\cosh u}\dd u\)，于是
\begin{equation}
\frac{\dd^2\Gamma_{\rm EELS}}{\dd z\,\dd\omega}
=\frac{q^2}{2\pi^2\varepsilon_0\hbar v^2}
K_0\!\left(\frac{2\omega b}{v}\right)
\bigl[-\operatorname{Im}r_{\rm es}(\omega)\bigr].
\label{eq:feqo-planar-eels-k0}
\end{equation}
左边单位为 \(\si{\second\per\metre}\)，与右边的 \(q^2/(\varepsilon_0\hbar v^2)\) 相同；乘 \(\dd z\dd\omega\) 后才是无量纲概率。\(K_0\) 随 \(\omega b/v\) 增大而衰减，表示离界面越远，高频近场越难被电子采到。

衰减还可以定量估计。对 \(x=2\omega b/v\gg1\)，积分 \(K_0(x)=\int_0^\infty e^{-x\cosh u}\dd u\) 主要来自 \(u\simeq0\)。在这里 \(\cosh u=1+u^2/2+O(u^4)\)，所以把积分的主贡献换成高斯积分给
\[
K_0(x)=\sqrt{\frac{\pi}{2x}}e^{-x}
\left[1+O(x^{-1})\right].
\]
因此把电子离界面距离增加 \(\Delta b\)，高频信号的主要指数因子减少 \(e^{-2\omega\Delta b/v}\)。这正是实验中高频近场对轨迹位置格外敏感的原因；若束斑宽度与 \(v/(2\omega)\) 可比，就应对不同 \(b\) 的概率平均，而不能只代入平均轨迹。

若取 Drude 材料 \(\varepsilon(\omega)=1-\omega_p^2/[\omega(\omega+\ii\gamma)]\)，弱阻尼 \(\gamma\ll\omega_p\) 时 \(1+\varepsilon\) 在 \(\omega\simeq\omega_p/\sqrt2\) 附近变小，损失谱出现表面模峰；\(\gamma\) 控制峰宽。无限平面沿 \(z\) 保持平移对称，电子给出的 \(k_z=\omega/v>\omega/c\) 不在真空传播光锥内，所以在此理想几何中远场阴极发光为零。材料吸收或表面波仍可让\cref{eq:feqo-planar-eels-k0} 为正。这是损失与发光不能混同的一个完整反例。

这条 Drude 谱的形状可以直接从\cref{eq:feqo-planar-eels-k0} 数值积出来。\cref{fig:planar-drude-eels} 取 \(\hbar\omega_p=9.0\,\si{\electronvolt}\)、\(\hbar\gamma=0.08\,\si{\electronvolt}\)、\(v=0.70c\)、\(b=20\,\si{\nano\metre}\)，实线是含推迟效应的完整 Weyl 积分，虚线标出无损极限下的同步极点。谱峰落在 \(5.25\,\si{\electronvolt}\)，几乎正对虚线给出的 \(5.23\,\si{\electronvolt}\)；阻尼把本该发散的极点削成有限高度 \(2.6\times10^{-4}\,\si{\per\electronvolt\per\nano\metre}\)，峰形明显偏向高能侧——\(6.0\,\si{\electronvolt}\) 处仍有峰值的四分之一，而低能侧 \(4.5\,\si{\electronvolt}\) 已不到峰值的百分之一。虚线与实线的重合说明摆位预测由 \(1+\varepsilon=0\) 给出，峰高和峰宽才需要 \(b\) 与 \(\gamma\) 入场。

\begin{figure}[htbp]
\centering
\includegraphics[width=0.86\linewidth]{planar_drude_eels.pdf}
\caption{真空侧掠射电子对 Drude 半空间的单位长度 EELS 谱。实线为完整 Weyl 积分数值结果，虚线为无损同步极点 \(\hbar\omega=\hbar\omega_p/\sqrt{1-\varepsilon_{\rm sync}}\simeq5.23\,\si{\electronvolt}\)；参数取 \(\hbar\omega_p=9.0\,\si{\electronvolt}\)、\(\hbar\gamma=0.08\,\si{\electronvolt}\)、\(\beta=0.70\)、\(b=20\,\si{\nano\metre}\)。纵轴单位为 \(\si{\per\electronvolt\per\nano\metre}\)。}
\label{fig:planar-drude-eels}
\end{figure}

峰宽的说法可由复介电函数而非图像推出。写 \(\omega_s=\omega_p/\sqrt2\)、\(\delta=\omega-\omega_s\)，并取 \(|\delta|,\gamma\ll\omega_s\)。逐项展开得
\[
1+\operatorname{Re}\varepsilon\simeq\frac{4\delta}{\omega_s},
\qquad
\operatorname{Im}\varepsilon\simeq\frac{2\gamma}{\omega_s},
\qquad
-\operatorname{Im}r_{\rm es}\simeq
\frac{\gamma\omega_s}{4\delta^2+\gamma^2}.
\]
在 \(K_0(2\omega b/v)\) 相对这个窄峰变化缓慢时，损失峰的全半高宽近似为 \(\gamma\)。若 \(b\) 很大使 \(K_0\) 在峰宽内显著变化，测得的峰形会被轨迹选择重新加权；若迟滞不可忽略，静电反射系数本身也必须换成 Maxwell 边界问题的结果。

\begin{exercise}
从界面两侧电势连续和 \(\varepsilon\partial_x\varphi\) 连续推导\cref{eq:feqo-planar-image-potential} 中的反射系数，并求无损 Drude 极限的峰频率。
\end{exercise}
